% Part: first-order-logic
% Chapter: lindstrom
% Section: introduction

\documentclass[../../../include/open-logic-section]{subfiles}

\begin{document}

\olfileid{mod}{lin}{int}
\section{تعارف}

ایس باب وچ ساڈا مقصد Lindstr\"om دی اوہ خصوصیت بندی ثابت کرنا اے جس
دے مطابق پہلے درجے دی منطق زیادہ توں زیادہ وسیع منطق اے جس لئی (کجھ ہور
شرطاں دے تحت) محدود گواہی دا قضیہ تے نزولی L\"owenheim--Skolem دا
قضیہ قائم رہندے نیں
(\olref[fol][com][com]{thm:compactness} تے
\olref[fol][com][dls]{thm:downward-ls})۔ پہلے سانوں منطقاں دے
اُس عمومی طبقے دی ہور وسیع خصوصیت بندی چاہیدی اے جس اُتے ایہ
قضیہ لاگو ہوندا اے۔ اسیں اپنے آپ نوں \emph{تعلقی} زباناں تک
محدود رکھاں گے، یعنی ایہو جیہیاں زباناں جنہاں وچ صرف
!!{predicate}s تے انفرادی مستقلاں ہوندے نیں، پر !!{function}s
نہیں ہوندے۔

\end{document}
